% =============================================================================
%  Front matter: Lời cam đoan, Lời cảm ơn, Tóm tắt, Mục lục, Danh mục
% =============================================================================

% --------- LỜI CAM ĐOAN -----------------------------------------------------
\clearpage
\thispagestyle{empty}
\begin{center}
{\fontsize{14}{17}\bfseries\MakeUppercase Lời cam đoan\par}
\end{center}
\vspace{1cm}

\hspace{1.27cm}Tôi xin cam đoan đây là công trình nghiên cứu của riêng tôi. Các
số liệu, kết quả thực nghiệm và kết luận nêu trong Luận văn là trung thực và
chưa từng được công bố trong bất kỳ công trình nào khác. Toàn bộ mã nguồn,
dữ liệu huấn luyện, các hình ảnh và bảng biểu trình bày trong Luận văn đều
do tôi tự xây dựng, ngoại trừ những phần đã trích dẫn nguồn rõ ràng trong
mục Tài liệu tham khảo theo chuẩn IEEE.

\vspace{0.5cm}
\hspace{1.27cm}Tôi xin chịu trách nhiệm về đề tài nghiên cứu của mình.

\vspace{2cm}
\begin{flushright}
\begin{minipage}{0.45\textwidth}
\centering
TP. Hồ Chí Minh, ngày \ldots\ tháng \ldots\ năm 2026\\[0.3cm]
Tác giả Luận văn\\[2.5cm]
\textbf{Nguyễn Duy Thanh}
\end{minipage}
\end{flushright}

% --------- LỜI CẢM ƠN -------------------------------------------------------
\clearpage
\thispagestyle{empty}
\begin{center}
{\fontsize{14}{17}\bfseries\MakeUppercase Lời cảm ơn\par}
\end{center}
\vspace{1cm}

\hspace{1.27cm}Để hoàn thành Luận văn Thạc sĩ này, tôi đã nhận được sự giúp đỡ
và hỗ trợ quý báu từ nhiều cá nhân và tổ chức. Tôi xin được bày tỏ lòng biết ơn
sâu sắc đến:

\vspace{0.3cm}
\hspace{1.27cm}\textbf{TS. Trần Thiên Thanh} và \textbf{TS. Ngô Hoàng Tú},
hai người hướng dẫn khoa học, đã tận tình chỉ bảo, định hướng nghiên cứu
và động viên tôi trong suốt quá trình thực hiện đề tài. Những góp ý quý
báu của hai Thầy về cả mặt lý thuyết và triển khai mô phỏng đã giúp Luận
văn có được kết quả như hiện tại.

\vspace{0.3cm}
\hspace{1.27cm}\textbf{Quý Thầy, Cô Viện Đào tạo Sau đại học và Khoa Công
nghệ Thông tin, Trường Đại học Giao thông vận tải TP. Hồ Chí Minh} đã
trang bị cho tôi nền tảng kiến thức vững chắc về khoa học máy tính, trí
tuệ nhân tạo và xử lý dữ liệu, là cơ sở thiết yếu để thực hiện công trình
này.

\vspace{0.3cm}
\hspace{1.27cm}\textbf{Cộng đồng nghiên cứu mã nguồn mở} (PyTorch, NumPy,
Matplotlib, draw.io) đã cung cấp các công cụ giúp việc cài đặt, huấn luyện
và trực quan hóa kết quả thực nghiệm trở nên thuận tiện và đáng tin cậy.

\vspace{0.3cm}
\hspace{1.27cm}\textbf{Gia đình và bạn bè} đã luôn ở bên động viên, chia sẻ và
hỗ trợ tôi cả về vật chất lẫn tinh thần trong suốt thời gian học tập và
nghiên cứu.

\vspace{0.5cm}
\hspace{1.27cm}Mặc dù đã có nhiều cố gắng, song do giới hạn về thời gian và
kinh nghiệm, Luận văn không tránh khỏi những thiếu sót. Tôi rất mong nhận
được ý kiến đóng góp của quý Thầy/Cô và bạn đọc để Luận văn được hoàn thiện
hơn.

\vspace{0.5cm}
\hspace{1.27cm}Xin trân trọng cảm ơn!

\vspace{1.5cm}
\begin{flushright}
\begin{minipage}{0.45\textwidth}
\centering
TP. Hồ Chí Minh, ngày \ldots\ tháng \ldots\ năm 2026\\[0.3cm]
Tác giả Luận văn\\[2.5cm]
\textbf{Nguyễn Duy Thanh}
\end{minipage}
\end{flushright}

% --------- MỤC LỤC ----------------------------------------------------------
\clearpage
\thispagestyle{empty}
\renewcommand{\contentsname}{\hfill MỤC LỤC \hfill}
\tableofcontents

% --------- DANH MỤC TỪ VIẾT TẮT ---------------------------------------------
\clearpage
\thispagestyle{empty}
\begin{center}
{\fontsize{14}{17}\bfseries\MakeUppercase Danh mục các ký hiệu, các chữ viết tắt\par}
\end{center}
\vspace{0.4cm}

\noindent\textbf{1. Các ký hiệu toán học}\par
\vspace{0.2cm}

\renewcommand{\arraystretch}{1.2}
\setlength{\tabcolsep}{4pt}
\begin{longtable}{p{3.5cm} p{11cm}}
\toprule
\textbf{Ký hiệu} & \textbf{Ý nghĩa} \\
\midrule
\endhead
$N$ & Số phần tử của STAR-RIS \\
$K$ & Tổng số người dùng (UE) trong hệ thống \\
$K_R, K_T$ & Số UE trong vùng phản xạ R và vùng truyền qua T \\
$M$ & Số anten của trạm gốc (BS) \\
$h_{d,k}$ & Kênh truyền trực tiếp từ BS đến UE thứ $k$ \\
$\mathbf{G}$ & Vector kênh truyền từ BS đến STAR-RIS, $\mathbf{G} \in \mathbb{C}^{N \times 1}$ \\
$\mathbf{h}_{r,k}^X$ & Vector kênh từ STAR-RIS đến UE thứ $k$, vùng $X \in \{R, T\}$ \\
$h_{eff,k}$ & Kênh hiệu dụng SISO tương đương BS--UE thứ $k$ \\
$\Phi_n^X$ & Hệ số phản xạ/truyền qua của phần tử thứ $n$, $X \in \{R, T\}$ \\
$\boldsymbol{\Phi}_X$ & Ma trận pha dịch STAR-RIS vùng $X$, dạng đường chéo \\
$\beta_n^R, \beta_n^T$ & Hệ số phân chia năng lượng phản xạ/truyền qua, $\in [0,1]$ \\
$\theta_n^R, \theta_n^T$ & Pha phản xạ/truyền qua của phần tử thứ $n$, $\in [0, 2\pi)$ \\
$\sqrt{\beta_n^X}$ & Biên độ phức của hệ số STAR-RIS \\
$x$ & Tín hiệu phát hợp nhất tại BS \\
$y_k$ & Tín hiệu thu tại UE thứ $k$ \\
$s_c, s_k$ & Luồng chung và luồng riêng (mã hóa) của UE thứ $k$ \\
$W_k, W_k^c, W_k^p$ & Thông điệp gốc, phần chung, phần riêng của UE thứ $k$ \\
$P_c, P_k$ & Công suất phát dành cho luồng chung và luồng riêng UE thứ $k$ \\
$P_{\max}$ & Công suất phát tối đa của BS \\
$c_k$ & Tỷ lệ phân chia luồng chung cho UE thứ $k$, $\sum_k c_k = 1$ \\
$\sigma^2$ & Công suất tạp âm Gauss phức tại UE \\
$\gamma_{c,k}, \gamma_{p,k}$ & SINR luồng chung và luồng riêng tại UE thứ $k$ \\
$R_{c,k}, R_{p,k}$ & Tốc độ luồng chung khả thi và luồng riêng tại UE thứ $k$ \\
$R_c$ & Tốc độ luồng chung của toàn hệ thống, $R_c = \min_k R_{c,k}$ \\
$R_k$ & Tổng tốc độ của UE thứ $k$, $R_k = c_k R_c + R_{p,k}$ \\
$R_{sum}$ & Tổng tốc độ toàn hệ thống, $R_{sum} = \sum_k R_k$ \\
$R_{\min}$ & Yêu cầu QoS tối thiểu cho từng UE \\
$\alpha_d, \alpha_G, \alpha_r$ & Số mũ suy hao đường truyền cho từng loại kênh \\
$d_0, \mathrm{PL}_0$ & Khoảng cách tham chiếu và suy hao tham chiếu \\
$\mathcal{S}, \mathcal{A}, \mathcal{R}$ & Không gian trạng thái, hành động, phần thưởng (MDP) \\
$s_t, a_t, r_t$ & Trạng thái, hành động, phần thưởng tại bước thời gian $t$ \\
$o_i, a_i$ & Quan sát cục bộ và hành động của tác tử thứ $i$ \\
$\mu_i, Q_i$ & Mạng Actor và Critic của tác tử thứ $i$ \\
$\theta_{\mu_i}, \phi_i$ & Tham số trọng số mạng Actor và Critic \\
$\gamma$ & Hệ số chiết khấu của MDP, $\gamma \in [0,1)$ \\
$\tau$ & Hệ số cập nhật mềm cho mạng mục tiêu \\
$\lambda_t$ & Hệ số phạt Lagrangian thích nghi tại bước $t$ \\
$\Delta\phi_n^X$ & Phần dư pha (residual) đầu ra của tác tử pha \\
$\mathcal{N}_t^{\text{OU}}$ & Nhiễu khám phá Ornstein--Uhlenbeck \\
$|\mathcal{D}|$ & Kích thước bộ đệm phát lại (replay buffer) \\
$\mathds{1}[\cdot]$ & Hàm chỉ thị (indicator function) \\
\bottomrule
\end{longtable}

\vspace{0.4cm}
\noindent\textbf{2. Các chữ viết tắt}\par
\vspace{0.2cm}

\renewcommand{\arraystretch}{1.2}
\setlength{\tabcolsep}{4pt}
\begin{longtable}{p{2cm} p{7cm} p{5.5cm}}
\toprule
\textbf{Viết tắt} & \textbf{Tiếng Anh} & \textbf{Tiếng Việt} \\
\midrule
\endhead
6G   & Sixth Generation                       & Thế hệ thứ sáu \\
AI   & Artificial Intelligence                & Trí tuệ nhân tạo \\
BCD  & Block Coordinate Descent               & Hạ dần khối tọa độ \\
BS   & Base Station                           & Trạm gốc \\
CSI  & Channel State Information              & Thông tin trạng thái kênh \\
CTDE & Centralised Training, Decentralised Execution & Huấn luyện tập trung, thực thi phân tán \\
DDPG & Deep Deterministic Policy Gradient     & Gradient chính sách tất định sâu \\
DL   & Deep Learning                          & Học sâu \\
DRL  & Deep Reinforcement Learning            & Học tăng cường sâu \\
DQN  & Deep Q-Network                         & Mạng Q sâu \\
ES   & Energy Splitting                       & Chia tách năng lượng \\
ISAC & Integrated Sensing and Communication   & Cảm biến và truyền thông tích hợp \\
LOS  & Line of Sight                          & Đường truyền thẳng \\
MA   & Multiple Access                        & Đa truy nhập \\
MADDPG & Multi-Agent DDPG                     & DDPG đa tác tử \\
MARL & Multi-Agent Reinforcement Learning     & Học tăng cường đa tác tử \\
MDP  & Markov Decision Process                & Quá trình ra quyết định Markov \\
MIMO & Multiple-Input Multiple-Output         & Nhiều ngõ vào nhiều ngõ ra \\
MLP  & Multi-Layer Perceptron                 & Mạng nơ-ron đa lớp \\
MS   & Mode Switching                         & Chuyển đổi chế độ \\
NLOS & Non-Line-of-Sight                      & Không có đường truyền thẳng \\
NN   & Neural Network                         & Mạng nơ-ron \\
NOMA & Non-Orthogonal Multiple Access         & Đa truy nhập không trực giao \\
OFDMA & Orthogonal Frequency-Division MA      & Đa truy nhập phân chia theo tần số trực giao \\
OMA  & Orthogonal Multiple Access             & Đa truy nhập trực giao \\
OU   & Ornstein--Uhlenbeck                    & Nhiễu Ornstein--Uhlenbeck \\
PPO  & Proximal Policy Optimisation           & Tối ưu hóa chính sách lân cận \\
QoS  & Quality of Service                     & Chất lượng dịch vụ \\
ReLU & Rectified Linear Unit                  & Đơn vị tuyến tính chỉnh lưu \\
RF   & Radio Frequency                        & Tần số vô tuyến \\
RIS  & Reconfigurable Intelligent Surface     & Bề mặt phản xạ thông minh \\
RL   & Reinforcement Learning                 & Học tăng cường \\
RSMA & Rate-Splitting Multiple Access         & Đa truy nhập theo phân tách tốc độ \\
SDMA & Space-Division Multiple Access         & Đa truy nhập phân chia theo không gian \\
SIC  & Successive Interference Cancellation   & Khử nhiễu liên tiếp \\
SINR & Signal-to-Interference-plus-Noise Ratio & Tỷ số tín hiệu trên nhiễu cộng tạp âm \\
SISO & Single-Input Single-Output             & Một ngõ vào một ngõ ra \\
STAR-RIS & Simultaneously Transmitting And Reflecting RIS & RIS truyền-phản xạ đồng thời \\
SWIPT & Simultaneous Wireless Information and Power Transfer & Truyền tin và điện đồng thời \\
TD3  & Twin Delayed DDPG                      & DDPG sinh đôi trì hoãn \\
TS   & Time Switching                         & Chuyển đổi thời gian \\
UAV  & Unmanned Aerial Vehicle                & Phương tiện bay không người lái \\
UE   & User Equipment                         & Thiết bị người dùng \\
URLLC & Ultra-Reliable Low-Latency Comms.     & Truyền thông siêu tin cậy độ trễ thấp \\
\bottomrule
\end{longtable}

% --------- DANH MỤC BẢNG ----------------------------------------------------
\clearpage
\thispagestyle{empty}
\renewcommand{\listtablename}{\hfill DANH MỤC CÁC BẢNG \hfill}
\listoftables

% --------- DANH MỤC HÌNH ----------------------------------------------------
\clearpage
\thispagestyle{empty}
\renewcommand{\listfigurename}{\hfill DANH MỤC CÁC HÌNH ẢNH \hfill}
\listoffigures
